% !TeX program = pdflatex
% =============================================================================
% Prüfungsreglement Physik – Kantonsschule KFR
%
% Allgemeines, schülerinnen- und schülergerichtetes Informationsblatt: Wie
% ist eine Physikprüfung aufgebaut, welche Regeln gelten bei der Korrektur
% und Bewertung, was ist mitzubringen. Nicht themengebunden, gilt für jede
% Prüfung, sofern eine einzelne Prüfung nicht ausdrücklich etwas anderes
% festlegt (z. B. eine abweichende Punktzahl G für Note 6).
% =============================================================================
\documentclass[11pt,a4paper]{article}

\usepackage[T1]{fontenc}
\usepackage[utf8]{inputenc}
\usepackage[ngerman]{babel}
\usepackage{lmodern}
\usepackage[a4paper,margin=2.2cm,headheight=14pt]{geometry}
\usepackage{amsmath}
\usepackage{siunitx}
% Hausstil: zusammengesetzte Einheiten immer als echter Bruch (\frac), nicht
% als Inline-Schrägstrich oder \tfrac -- gilt global für jedes \si/\SI.
\sisetup{per-mode=fraction,fraction-command=\frac}
\usepackage{array,booktabs,tabularx}
\usepackage{enumitem}
\usepackage{xcolor}
\usepackage[most]{tcolorbox}
\usepackage{lastpage}
\usepackage{fancyhdr}
\usepackage{setspace}

\setlength{\parindent}{0pt}
\setlength{\parskip}{0.6em}
\setstretch{1.08}
\setlist[itemize]{itemsep=0.4em}

% -----------------------------
% Farben (Corporate-Farbschema Physik KFR, wie in den übrigen Vorlagen)
% -----------------------------
\definecolor{titleblue}{HTML}{183A6B}
\definecolor{accentblue}{HTML}{2E6F9E}
\definecolor{lightblue}{HTML}{EAF4FB}
\definecolor{verylightblue}{HTML}{F7FBFE}
\definecolor{darkgray}{HTML}{4A4A4A}

\newcolumntype{Y}{>{\centering\arraybackslash}X}

% -----------------------------
% Boxen
% -----------------------------
\tcbset{before skip=10pt, after skip=10pt}
\newtcolorbox{titlebox}{
  enhanced,
  colback=titleblue,
  colframe=titleblue,
  boxrule=0pt,
  arc=3mm,
  left=3mm,right=3mm,top=3mm,bottom=3mm
}

\newtcolorbox{infobox}[1][]{
  enhanced,
  breakable,
  colback=lightblue,
  colframe=accentblue,
  title={#1},
  fonttitle=\bfseries,
  boxrule=0.7pt,
  arc=2mm,
  left=6mm,right=6mm,top=3.5mm,bottom=3.5mm
}

\newtcolorbox{beispielbox}[1][Beispiel]{
  enhanced,
  breakable,
  colback=verylightblue,
  colframe=accentblue,
  title={#1},
  fonttitle=\bfseries,
  boxrule=0.7pt,
  arc=2mm,
  left=6mm,right=6mm,top=3.5mm,bottom=3.5mm
}

% --- Kopfzeile ---------------------------------------------------------------
\newcommand{\Kopfzeile}[4]{%
  \begin{titlebox}
    {\color{white}\sffamily\bfseries\Large #4}
  \end{titlebox}
  \noindent\textbf{Klasse:}~#1 \hfill \textbf{Datum:}~#2 \hfill \textbf{Thema:}~#3
}

\pagestyle{fancy}
\fancyhf{}
\rhead{\textcolor{darkgray}{Prüfungsreglement Physik}}
\lhead{\textcolor{darkgray}{Kantonsschule KFR}}
\cfoot{\textcolor{darkgray}{Seite \thepage\ von \pageref{LastPage}}}
\renewcommand{\headrulewidth}{0.4pt}

\begin{document}

\Kopfzeile{alle Klassen}{\today}{Prüfungen}{Prüfungsreglement: Aufbau und Regeln}

Dieses Blatt erklärt dir, wie eine Physikprüfung bei mir aufgebaut ist und
nach welchen Regeln sie korrigiert und bewertet wird. Es gilt für jede
Physikprüfung, sofern auf der Prüfung selbst nichts anderes vermerkt ist.

\begin{infobox}[Aufbau der Prüfung]
\begin{itemize}
  \item \textbf{Erste Seite:} Nach der Kopfzeile mit Klasse, Datum, Thema
    und Titel folgt das Bewertungsraster. Es zeigt dir, wie viele Punkte
    jede Aufgabe maximal gibt, und dient vor allem der Korrektur als
    Übersicht.
  \item \textbf{Formelsammlung:} Zu jeder Prüfung gehört eine
    Formelsammlung mit allen relevanten Formeln. Sie ist entweder direkt
    in der Prüfung abgedruckt oder liegt als separates Blatt bei.
  \item \textbf{Themenabdeckung:} Zu jedem im Unterricht behandelten Thema
    gibt es mindestens eine Aufgabe.
  \item \textbf{Aufgabentypen:} Die Prüfung deckt verschiedene Fähigkeiten
    ab. Es gibt Verständnisaufgaben (Multiple Choice oder Textaufgaben mit
    Erklärung) und Rechenaufgaben.
  \item \textbf{Schwierigkeitsaufbau:} Die ersten Aufgaben sind einfach,
    danach folgen mehrere mittelschwere Aufgaben, und die letzten Aufgaben
    sind anspruchsvoll.
  \item \textbf{Platz zum Lösen:} Du löst die Aufgaben direkt auf dem
    dafür vorgesehenen Platz in der Prüfung.
  \item \textbf{Letzte Seite:} Hier stehen deine Note aus der Prüfung und,
    daneben, die mündliche Note für den Unterricht.
\end{itemize}
\end{infobox}

\section*{Bewertung: von Punkten zur Note}

Das Bewertungsschema ist linear: Jeder zusätzliche Punkt bringt gleich
viel Note wie jeder andere. Die Höchstnote musst du nicht zwingend mit der
vollen Maximalpunktzahl erreichen, die dafür nötige Punktzahl ist auf dem
Bewertungsraster der jeweiligen Prüfung angegeben.

\section*{Genauigkeit der Resultate}

Das Resultat einer Rechenaufgabe wird mit derselben Genauigkeit
(Anzahl signifikanter Stellen) angegeben wie die ungenaueste in der
Aufgabe gegebene Grösse. Weder mehr Nachkommastellen (täuscht eine
Genauigkeit vor, die die Rechnung gar nicht hergibt) noch weniger
(verschenkt vorhandene Information) sind korrekt.

\begin{beispielbox}[Beispiel: Genauigkeit]
Gegeben: $m = \SI{2.5}{\kilo\gram}$ (2 signifikante Stellen),
$v = \SI{3.10}{\meter\per\second}$ (3 signifikante Stellen).

Gesucht: kinetische Energie $E = \frac{1}{2} m v^2$.

\[
E = \frac{1}{2} \cdot \SI{2.5}{\kilo\gram} \cdot (\SI{3.10}{\meter\per\second})^2
  = \SI{12.0125}{\joule}
\]

Da $m$ mit nur 2 signifikanten Stellen gegeben ist, wird auch das
Resultat auf 2 signifikante Stellen gerundet: $E \approx \SI{12}{\joule}$.
Weder $\SI{12.0125}{\joule}$ (zu genau) noch $\SI{10}{\joule}$ (zu
ungenau gerundet) wären korrekt.
\end{beispielbox}

\section*{Der Lösungsweg: so entstehen die Punkte einer Rechenaufgabe}

Der grösste Teil der Punkte einer Rechenaufgabe kommt vom Lösungsweg,
nicht vom blossen Endresultat. Der Lösungsweg muss deshalb immer
vollständig aufgeschrieben werden, in dieser Reihenfolge:

\begin{enumerate}
  \item die verwendete Formel (Ansatz),
  \item die Formel nach der gesuchten Grösse aufgelöst, symbolisch, also
    noch ohne eingesetzte Zahlenwerte,
  \item erst danach die gegebenen Werte eingesetzt.
\end{enumerate}

Die Maximalpunktzahl einer Rechenaufgabe setzt sich damit so zusammen:

\[
P_{\text{Aufgabe}} = P_{\text{Formel}} + P_{\text{Umstellen}}
  + P_{\text{Einsetzen}} + P_{\text{Resultat}} + P_{\text{Einheit}}
\]

Das richtige Resultat (Zahlenwert) gibt den vorletzten Punkt, die
richtige Einheit den letzten. Fehlt die Einheit beim Resultat, geht
dieser letzte Punkt verloren, auch wenn der Zahlenwert stimmt.

\begin{beispielbox}[Beispiel: Punkteschema]
Ein Auto beschleunigt aus dem Stillstand mit
$a = \SI{4.0}{\meter\per\second\squared}$ auf $v = \SI{20}{\meter\per\second}$.
Berechnen Sie die dafür benötigte Zeit $t$. [5 P.]

\begin{tabularx}{\textwidth}{@{}X c@{}}
\toprule
Schritt & Punkte \\
\midrule
Formel (Ansatz): $v = v_0 + a t$ & 1 \\
Auflösen nach $t$: $t = \dfrac{v - v_0}{a}$ & 1 \\
Einsetzen: $t = \dfrac{\SI{20}{\meter\per\second} - 0}{\SI{4.0}{\meter\per\second\squared}}$ & 1 \\
Resultat: $t = 5.0$ & 1 \\
Einheit: $\si{\second}$ & 1 \\
\midrule
Total & 5 \\
\bottomrule
\end{tabularx}
\end{beispielbox}

\section*{Unabhängige Teilaufgaben und Folgefehler}

Teilaufgaben werden nach Möglichkeit so gestellt, dass sie unabhängig
voneinander lösbar sind. Baut eine Teilaufgabe zwingend auf dem Resultat
einer vorherigen auf, wird in der Aufgabenstellung ein Ersatzwert
angegeben, mit dem weitergerechnet werden kann, falls die vorherige
Teilaufgabe nicht gelöst wurde. Ein Rechenfehler in einer frühen
Teilaufgabe wird nur einmal als Fehler gewertet: Wird das eigene, falsche
Zwischenresultat in einer späteren Teilaufgabe mit richtigem Ansatz
korrekt weiterverwendet, gibt es dafür die volle Punktzahl.

\begin{beispielbox}[Beispiel: Folgefehler]
Ein Ball wird mit $v_0 = \SI{15}{\meter\per\second}$ senkrecht nach oben
geworfen ($g = \SI{9.81}{\meter\per\second\squared}$).

\textbf{a) [2 P.]} Berechnen Sie die Zeit $t_1$ bis zum höchsten Punkt
($v = 0$).
\[
t_1 = \frac{v_0}{g} = \frac{\SI{15}{\meter\per\second}}{\SI{9.81}{\meter\per\second\squared}} \approx \SI{1.5}{\second}
\]

\textbf{b) [3 P.]} Berechnen Sie mit $t_1$ die maximale Höhe $h$ über dem
Abwurfpunkt. Haben Sie a) nicht gelöst, verwenden Sie
$t_1 = \SI{2.0}{\second}$.
\[
h = v_0 t_1 - \frac{1}{2} g t_1^2
  = \SI{15}{\meter\per\second} \cdot \SI{1.5}{\second}
  - \frac{1}{2} \cdot \SI{9.81}{\meter\per\second\squared} \cdot (\SI{1.5}{\second})^2
  \approx \SI{11}{\meter}
\]

Hätte eine Schülerin in a) fälschlicherweise $t_1 = \SI{2.0}{\second}$
berechnet (statt der korrekten $\SI{1.5}{\second}$) und diesen Wert in b)
mit dem richtigen Ansatz weiterverwendet
($h \approx \SI{10}{\meter}$), erhielte sie in b) trotzdem die volle
Punktzahl. Der Fehler wird nur in a) bestraft, nicht ein zweites Mal in b).
\end{beispielbox}

\section*{Erlaubte Hilfsmittel}

\begin{infobox}[Mitzubringen]
\begin{itemize}
  \item ein nicht radierbarer Stift in blauer oder schwarzer Farbe
    (z.\,B. Kugelschreiber),
  \item ein von der Schule zugelassener Taschenrechner,
  \item ein Lineal,
  \item eine Formelsammlung, sofern diese nicht Teil der Prüfung ist.
\end{itemize}
\end{infobox}

\begin{infobox}[Zusätzliche Blätter]
Reicht der zur Verfügung gestellte Platz nicht aus, können vor der
Prüfung zusätzliche leere Blätter bei der Lehrperson verlangt werden.
Jedes zusätzliche Blatt muss mit Name und der zugehörigen Aufgabennummer
beschriftet werden. Unbeschriftete Zusatzblätter werden nicht bewertet.
\end{infobox}

\end{document}
